\documentclass[11pt]{letter}
\usepackage[margin=2.5cm]{geometry}
\signature{First Author\\ on behalf of all authors}
\address{Department, Institution\\ City, Country\\ first.author@institution.edu}
\begin{document}
\begin{letter}{The Editor-in-Chief\\ Measurement}
\opening{Dear Editor,}

We submit the manuscript \emph{``Do vibration-based fault classifiers generalise to unseen rotation speeds? A
leakage-free benchmark of feature models, Kolmogorov--Arnold, Hamiltonian and order-conditioned Fourier neural operator
networks on MAFAULDA and COMFAULDA''} for consideration as a research article in \emph{Measurement}.

Accuracies close to 99\,\% are routinely reported on public machinery fault databases, but they are obtained with test
recordings made at speeds already present in training. We measure, under splits that never share a recording and that
hold out whole speed ranges, how much of that performance survives. Every model loses 0.11--0.32 macro-F1; a random
forest on order-domain features remains the most accurate and stable model, an order-conditioned Fourier neural
operator generalises clearly better for unbalance, and an ablation explains why the same speed conditioning hurts the
other faults. We also show that compound faults are not recognised from single-fault training, and we document
recording-session artefacts of MAFAULDA (tachometer faults, a drifting 1X phase reference, an anomalous accelerometer)
that matter for any method using phase or channel-level information.

The work concerns the measurement and evaluation of condition-monitoring methods, which we believe fits the scope of
\emph{Measurement}. All hypotheses and decision criteria were written before the experiments, results that failed
their criterion are reported, and the code, protocols and result files are openly available.

The manuscript is original, has not been published and is not under consideration elsewhere. All authors approved the
submission and declare no competing interests.

\closing{Yours sincerely,}
\end{letter}
\end{document}
